\begin{rubric}{Top 5 Most Relevant Publications (sorted by date)}

\entry*[2026] \textbf{Dependency-aware model repair: prioritizing repairs through consistency rule and inconsistency analysis} (\emph{Software and Systems Modeling}). 
This is the core contribution of my PhD thesis. It introduces a novel dependency-aware approach for model repair that prioritizes fixes based on their global impact, enabling more effective consistency maintenance. Demonstrates strong empirical results and extends our prior award-winning work (SANER 2024 Distinguished Paper Award).

\entry*[2024] \textbf{Contemporary Software Modernization: Strategies, Driving Forces, and Research Opportunities} (\emph{ACM Trans. Softw. Eng. Methodol.}). 
This SLR paper presents a comprehensive synthesis of the state of the art in software modernization, based on 126 studies. It identifies key strategies, industrial drivers, and research gaps, providing a widely applicable research agenda for both academia and practitioners.

\entry*[2023] \textbf{Fulfilling Industrial Needs for Consistency Among Engineering Artifacts} (\emph{ICSE-SEIP}). 
In this paper, we demonstrate a real-world applicability of my PhD research in collaboration with industry. It provides a scalable solution for maintaining consistency across heterogeneous engineering artifacts, validated on large industrial systems (20k+ elements) with real-time feedback.

\entry*[2022] \textbf{Software product line scoping: A systematic literature review} (\emph{Journal of Systems and Software}). 
This SLR establishes a consolidated view of SPL scoping by analyzing 58 studies and 41 approaches. It provides a generic scoping process and is frequently used as a reference for defining SPL scoping phases and practices.

\entry*[2020] \textbf{Systematic mapping study on domain-specific language development tools} (\emph{Empirical Software Engineering}). 
This large-scale SMS (230 papers, 59 tools) maps the landscape of DSL development environments, including methods, tools, and empirical evaluations. It is highly cited (100+ citations) and widely used as a reference for decision-making in DSL engineering and beyond.

\end{rubric}